(1) 相邻缝的相位差为 q，远场复振幅几何级数 \(A=A_1\sum_{j=0}^{N-1}e^{ijq}=A_1e^{i(N-1)q/2}\sin(Nq/2)/\sin(q/2)\)。故
\[
I/I_{\rm single}=[\sin(Nq/2)/\sin(q/2)]^2.
\]
在 \(q=2\pi m\) 处分子和分母同时为零，连续极限为 \(N^2\)，即相干主极大。

(2) 主峰满足 \(d\sin\theta_m=m\lambda\)。在其附近写 \(q=2\pi m+\delta\)，第一个零点来自 \(N\delta/2=\pm\pi\)，故相邻首个零点相对峰位 \(|\delta|=2\pi/N\)。因此 \(|\Delta(\sin\theta)|=\lambda/(Nd)\)，小角变化为
\[
\Delta\theta_{\rm zero}\simeq\frac{\lambda}{Nd\cos\theta_m}.
\]

(3) 固定 m,d，对光栅方程微分 \(d\cos\theta_m\,d\theta=m\,d\lambda\)，得
\[
\frac{d\theta_m}{d\lambda}=\frac{m}{d\cos\theta_m}.
\]
角散度随衍射级次增大，但必须满足 \(|m\lambda|\le d\)。

(4) 两近邻谱线峰位置相差约 \(|\delta\theta|=|d\theta/d\lambda|\,\delta\lambda\)。Rayleigh 判据令它等于单峰到首个零点的角间距，故 \(m\delta\lambda/(d\cos\theta)=\lambda/(Nd\cos\theta)\)，整理得 \(R=\lambda/\delta\lambda=mN\)。只有 m 非零、谱峰存在时此式有物理意义。
